% Comando de tab alineado con el formato APA (0.5 in)
\newcommand{\tab}{\hspace{1.27cm}}

% Comando para no tener que escribir el path de las imagenes cada vez, solo el nombre del archivo con su extensión.
\graphicspath{{assets/jpg/}{assets/png/}{assets/svg/}{assets/}}
\providecommand{\nat}[1]{#1}

%
\newcommand{\inicializarSeccion}[1]{
    % Use the document's section formatting (small-caps, centered)
    \section*{#1}
    \phantomsection
    \addcontentsline{toc}{subsection}{#1}
}

\newcommand{\inicializarSubseccion}[1]{
    % Mimic \inicializarSeccion but a bit smaller: use section* styling
    % while reducing the font size so it appears subordinate.
    \section*{\normalfont\normalsize\scshape\centering #1}
    \phantomsection
    \addcontentsline{toc}{subsubsection}{#1}
}

\newcommand{\figref}[1]{Figura~\ref{#1}}

\renewcommand{\figureautorefname}{Figura}
\renewcommand{\tableautorefname}{Tabla}
\renewcommand{\sectionautorefname}{Anexo}

\newcounter{anexo}
\setcounter{anexo}{0}

\makeatletter
\newcommand{\anexofigurecaption}[1]{%
    \refstepcounter{anexo}% incrementa y establece \@currentlabel provisional
    \protected@xdef\@currentlabel{Anexo \theanexo}% global
    \begingroup
    \captionsetup{labelformat=empty}% suprime "Figura N:" automático
    \captionof{figure}{\textbf{Anexo \theanexo:} #1}% texto de caption
    \endgroup
}

\newcommand{\anexotablecaption}[1]{%
    \refstepcounter{anexo}% incrementa y establece \@currentlabel provisional
    \protected@xdef\@currentlabel{Anexo \theanexo}% global
    \begingroup
    \captionsetup{labelformat=empty}% suprime "Tabla N:" automático
    \captionof{table}{\textbf{Anexo \theanexo:} #1}% texto de caption
    \endgroup
}
\makeatother
\renewcommand{\equationautorefname}{Ecuación}
\usepackage{cancel}

% ---------------------------------------------------------------------------
% \cite: el hipervínculo cubre la cita completa (autor, año) y apunta a la
% referencia bibliográfica, en lugar de enlazar solo el año (comportamiento
% por defecto de biblatex-apa).
% ---------------------------------------------------------------------------
\makeatletter

% Año de la cita sin enlace propio: el enlace lo pone el macro "cite" alrededor
% de toda la cita. Si el año se enlazara por su cuenta quedarían dos enlaces
% anidados sobre la misma referencia.
\newbibmacro*{cite:plabelyear+extradate:sin-enlace}{%
    \begingroup
        \let\bibhyperref\@firstofone
        \usebibmacro{cite:plabelyear+extradate}%
    \endgroup
}

% Igual que el macro "cite" de apa.cbx, pero con \bibhyperref (que enlaza a la
% referencia bibliográfica) alrededor de todo el contenido, de modo que autor,
% delimitador y año sean un solo enlace: [Naciones Unidas, 2015](enlace).
\renewbibmacro*{cite}{%
    \bibhyperref{%
        \iffieldequals{namehash}{\cbx@lasthash}
% Varias obras del mismo autor en un solo \cite
            {\setunit{\compcitedelim}%
                \usebibmacro{cite:plabelyear+extradate:sin-enlace}}%
% Cita normal
            {\ifnameundef{labelname}
% Sin autor/editor: se usa el título
                {\usebibmacro{cite:noname}%
                    \setunit{\printdelim{nameyeardelim}}%
                    \usebibmacro{cite:plabelyear+extradate:sin-enlace}%
                    \savefield{namehash}{\cbx@lasthash}}
% Con autor
                {\ifnameundef{shortauthor}
                    {\printnames{labelname}}%
                    {\cbx@apa@ifnamesaved
                        {\printnames{shortauthor}}
                        {\printnames[labelname]{author}%
                            \addspace\printnames[sabrackets]{shortauthor}}}%
                    \setunit{\printdelim{nameyeardelim}}%
                    \usebibmacro{cite:plabelyear+extradate:sin-enlace}%
                    \savefield{namehash}{\cbx@lasthash}}}%
    }%
    \setunit{\multicitedelim}%
}
\makeatother

% ---------------------------------------------------------------------------
% Pies de página. El estilo "fancy" replica lo que hacía \pagestyle{plain}
% (solo el número de página centrado, sin reglas); la nota de \pieDePagina
% no va en el pie sino al final del cuerpo del texto (ver más abajo).
% ---------------------------------------------------------------------------
\pagestyle{fancy}
\fancyhf{}
\fancyfoot[R]{\thepage\hspace{-2cm}}
\renewcommand{\headrulewidth}{0pt}
\renewcommand{\footrulewidth}{0pt}

% \pieDePagina{texto}: agrega "texto" como nota al pie de la página donde se
% escribe el comando, y solo de esa página. La nota se coloca al final del
% cuerpo del texto, precedida por una regla horizontal, y el número de página
% queda centrado debajo (en el margen inferior). El cuerpo de esa página se
% recorta el espacio que ocupe la nota, de modo que nunca se encima con el
% texto.
% El argumento puede ocupar varias líneas (con \\) y llevar formato, p. ej.:
%   \pieDePagina{\emph{Fecha:} Borrador para circulación, 17 de agosto de 2026.\\
%     \emph{Palabras clave:} problema, métodos, resultados.}
% Debe usarse en la página deseada, por ejemplo después de un \newpage o al
% comienzo de la sección que abre esa página.
\makeatletter
\long\def\pieDePagina#1{%
    \insert\footins{%
        \reset@font\footnotesize
        \interlinepenalty\interfootnotelinepenalty
        \splittopskip\footnotesep
        \splitmaxdepth\dp\strutbox
        \floatingpenalty\@MM
        \hsize\columnwidth
        \@parboxrestore
        \color@begingroup
            \rule\z@\footnotesep\ignorespaces#1\@finalstrut\strutbox
        \par
        \color@endgroup
    }%
}
\makeatother
